\section{Hyperbolic coordinates and the physical relative amplitude}
\label{sec:v14-normal-form}
The relative product has a local analytic explanation as well as the
half-line determinant construction. We give the explanation in physical
endpoint coordinates, where the observable is defined. The analytic
normal form recalled in De Simoi--Kaloshin--Leguil~\cite{DSKL}, printed
pp.~10 and 13 of arXiv:1905.00890v4, is a local statement; the symmetry
hypotheses of their global inverse theorem are not hypotheses for its
existence. The mixed-boundary and projection calculation below was
supplied in the author-requested referee benchmark of September 10,
2026. Its inclusion here makes the overlap of the two forward mechanisms
explicit. We do not attribute this calculation to the cited spectral
rigidity theorem.

\subsection{A relative estimate in mixed coordinates}
Let
\begin{equation}\label{eq:v14-normal-form}
 N(s,p)=(\Delta(sp)s,\Delta(sp)^{-1}p),
 \qquad \Delta(0)=\lambda\in(0,1),
\end{equation}
where $\Delta$ is real analytic and positive near zero. The integer $n$
counts iterates of this map. For the two-step return of the billiard,
$n$ counts returns, $j=2n$, and $\lambda=e^{-2\gamma}$.

\begin{lemma}[Normalized mixed derivative]
\label{lem:v14-mixed-normal-form}
There is a fixed neighborhood of $(s,t)=(0,0)$ on which the mixed
boundary problem $s_0=s$, $p_n=t$ has the unique small solution
\begin{equation}\label{eq:v14-mixed-equation}
 I=st\Delta(I)^n,\qquad
 p_0=t\Delta(I)^n,\qquad s_n=s\Delta(I)^n.
\end{equation}
All intermediate points stay in a fixed normal-form neighborhood. For
some $\rho\in(\lambda,1)$ and every fixed $k$,
\begin{equation}\label{eq:v14-twist-normal-form}
 T_n:=\left.\frac{\partial p_0}{\partial t}\right|_s
 =\frac{\Delta(I)^n}{1-nI\Delta'(I)/\Delta(I)},
 \qquad
 \|\lambda^{-n}T_n-1\|_{C^k}\le C_k\rho^n.
\end{equation}
These estimates hold for every $n\ge1$ on a common smaller neighborhood.
\end{lemma}
\begin{proof}
Choose $\lambda<\mu<\rho<1$ and a complex disk about zero on which
$\Delta$ is holomorphic, nonzero, and $|\Delta|\le\mu$.
Choose a complex endpoint bidisk $|s|,|t|<\epsilon$ so small that
$I\mapsto st\Delta(I)^n$ maps the chosen disk to itself and
\[
 \epsilon^2\|\Delta'\|_\infty
                   \sup_{n\ge1}n\mu^{n-1}<\tfrac12.
\]
The uniform contraction gives a holomorphic solution, with
$|I|\le\epsilon^2\mu^n$. The formulas for $p_0,s_n$ follow from
invariance of $sp$. At the intermediate index $i$, the coordinates are
$(s\Delta(I)^i,t\Delta(I)^{n-i})$, proving the locality assertion.

Put $D_n=1-nI\Delta'(I)/\Delta(I)$. Substitution of the fixed-point
identity shows that $|D_n|\ge1/2$. Differentiating the first identity in
\eqref{eq:v14-mixed-equation} gives $D_n I_t=s\Delta(I)^n$.
Differentiation of $p_0=t\Delta(I)^n$, without dividing by $s$, now
gives $T_n=\Delta(I)^n/D_n$. Thus the formula also holds at $s=0$.
Since the analytic logarithm of $\Delta(I)/\lambda$ vanishes at $I=0$,
\[
 \frac{\Delta(I)^n}{\lambda^n}
   =\exp\{n\log(\Delta(I)/\lambda)\}
   =1+O(n\mu^n),\qquad D_n=1+O(n\mu^n)
\]
uniformly on the complex bidisk. These are estimates of the
\emph{normalized} expressions. Cauchy estimates on a smaller fixed
bidisk give the same bound for every fixed number of endpoint
derivatives. Finally $n\mu^n\le C\rho^n$. No absolute error has
been divided by an uncontrolled exponentially small quantity.
\end{proof}

\subsection{Projection to physical endpoints}
Suppose the initial and final canonical charts are analytic maps
\[
 (u,P)=\mathcal C_L(s,p)=(U(s,p),R(s,p)),\qquad
 (v,Q)=\mathcal C_R(q,t)=(V(q,t),Z(q,t)),
\]
which send the origin to the origin and preserve the displayed
symplectic forms. Assume
\begin{equation}\label{eq:v14-transversality}
                 U_s(0,0)V_t(0,0)\ne0.
\end{equation}
For a physical generating function use the convention
$P=-W_{n,u}$, $Q=W_{n,v}$. Write $\alpha(u)$ and $\beta(v)$ for
the local inverses of $s\mapsto U(s,0)$ and $t\mapsto V(0,t)$.

\begin{proposition}[Physical projection and relative factorization]
\label{prop:v14-physical-normal-form}
There are a fixed physical endpoint box and $n_0$ such that, for
$n\ge n_0$, the mixed-to-physical endpoint map
\begin{equation}\label{eq:v14-endpoint-map}
 \Phi_n(s,t)=\bigl(U(s,t\Delta(I)^n),V(s\Delta(I)^n,t)\bigr)
\end{equation}
is invertible on that box. With $\mathcal D_n=\det D\Phi_n$, its
physical endpoint flux satisfies the exact identity
\begin{equation}\label{eq:v14-physical-jacobian}
 \mathcal J_n(u,v):=|W_{n,uv}(u,v)|
     =\left|\frac{T_n(s,t)}{\mathcal D_n(s,t)}\right|,
                 \qquad (s,t)=\Phi_n^{-1}(u,v).
\end{equation}
For every fixed $k$,
\begin{equation}\label{eq:v14-projected-limit}
 \left\|\frac{\mathcal J_n}{\mathcal J_n(0,0)}
                         -B_L\otimes B_R\right\|_{C^k}
                       \le C_k\rho^n,
\end{equation}
where the positive, individually normalized factors are
\begin{equation}\label{eq:v14-projection-amplitudes}
 B_L(u)=\frac{U_s(0,0)}{U_s(\alpha(u),0)},\qquad
 B_R(v)=\frac{V_t(0,0)}{V_t(0,\beta(v))}.
\end{equation}
Moreover, with $E_n=W_n-W_n(0,0)$,
\begin{equation}\label{eq:v14-action-limit}
 \|E_n-S_L\oplus S_R\|_{C^k}\le C_k\rho^n,
\end{equation}
where $S_L(0)=S_R(0)=0$ and
$S_L'(u)=-R(\alpha(u),0)$, $S_R'(v)=Z(0,\beta(v))$.
\end{proposition}
\begin{proof}
The preceding complex estimates, followed by Cauchy estimates, imply
\[
 \Phi_n=\Phi_\infty+O_{C^k}(\rho^n),\qquad
 \Phi_\infty(s,t)=(U(s,0),V(0,t)).
\]
Shrink the coordinate box so that both scalar projection derivatives
are bounded away from zero. Its limiting map is a diffeomorphism.
On a smaller physical box, the equation for $\Phi_n^{-1}$ is a
contraction about $\Phi_\infty^{-1}$ for all sufficiently large $n$.
Differentiating that equation yields convergence of the inverses in
every fixed derivative norm. In particular, $|\mathcal D_n|$ is bounded
away from zero on the relevant domain. This argument uses the same
physical box for all $n\ge n_0$.

Pull back the identity $\dd u\wedge\dd P=\dd s\wedge\dd p_0$
to the variables $(s,t)$. The right side is
$T_n\dd s\wedge\dd t$; the left side is
$(\partial P/\partial v)_u\mathcal D_n\dd s\wedge\dd t$.
This proves \eqref{eq:v14-physical-jacobian}. At the origin the exact
linear calculation gives
\begin{equation}\label{eq:v14-reference-jacobian}
 \mathcal D_n(0,0)=U_s(0,0)V_t(0,0)
             -\lambda^{2n}U_p(0,0)V_q(0,0),\qquad T_n(0,0)=\lambda^n.
\end{equation}
Also $\mathcal D_\infty(s,t)=U_s(s,0)V_t(0,t)$. Consequently the
normalized flux equals
\[
 \left|\frac{T_n(s,t)}{\lambda^n}
             \frac{\mathcal D_n(0,0)}{\mathcal D_n(s,t)}\right|,
 \qquad (s,t)=\Phi_n^{-1}(u,v).
\]
Apply \eqref{eq:v14-twist-normal-form}, convergence of the inverse maps,
and the nonzero determinant bound. The result is
\eqref{eq:v14-projected-limit}; the signs of the scalar derivatives
are constant, so their normalized ratios are positive.

Finally $p_0,s_n$ converge to zero in every fixed derivative norm.
The generating-function equations therefore give
\[
 W_{n,u}=-R(\alpha(u),0)+O_{C^k}(\rho^n),\qquad
 W_{n,v}= Z(0,\beta(v))+O_{C^k}(\rho^n).
\]
Integrate along endpoint segments, starting at the origin, to obtain
the zeroth-order estimate for $E_n$. The differentiated estimates give
all higher orders in \eqref{eq:v14-action-limit}.
\end{proof}

For the same physical section at both ends, $V=U$, formula
\eqref{eq:v14-reference-jacobian} gives
\[
 \mathcal J_n(0,0)=
 \frac{1}{|U_s(0,0)U_p(0,0)|}
                       \frac{\lambda^n}{1-\lambda^{2n}}.
\]
Thus the exact reference normalization has the hyperbolic-sine
correction, not merely its leading exponential. For the billiard
$j=2n$, the identity $\lambda=e^{-2\gamma}$ is consistent with
$d_{2n}^0=a_b/\sinh(2n\gamma)$.

\begin{remark}[A nonconstant amplitude with a linear normal form]
The projection dependence can be seen without a nonlinear $\Delta$.
Take $\Delta\equiv\lambda$ and the canonical charts
\[
 \mathcal C_L(s,p)=
  \left(s+p+as^2,\frac{-s+p+as^2}{2}\right),\qquad
 \mathcal C_R(q,t)=
  \left(q+t+bt^2,\frac{-q+t-bt^2}{2}\right).
\]
Their Jacobians equal one. The mixed endpoint determinant is
$(1+2as)(1+2bt)-\lambda^{2n}$, so the normalized flux, at
$(u,v)=\Phi_n(s,t)$, is exactly
\[
 \frac{1-\lambda^{2n}}
               {(1+2as)(1+2bt)-\lambda^{2n}}.
\]
The limiting physical factors are
$B_L(u)=(1+4au)^{-1/2}$ and $B_R(v)=(1+4bv)^{-1/2}$.
Both limiting action Hessians at zero equal $1/2$. This is a local
canonical example, not an asserted realization of arbitrary charts by
a billiard. It shows why linearity of the normal form does not force
constancy of its physical relative amplitude.
\end{remark}

\begin{corollary}[Identification of the analytic half-line amplitudes]
\label{cor:v14-amplitude-identification}
Apply Proposition~\ref{prop:v14-physical-normal-form} and
Theorem~\ref{thm:v4-factorization} to the same analytic, selected
even-flight billiard endpoint problem, in the same physical coordinates.
Then the functions in \eqref{eq:v14-projection-amplitudes} equal the
corresponding half-line determinant amplitudes. The actions $S_L,S_R$
equal the corresponding half-line actions, with their zero-value
normalizations. This identifies the analytic determinant amplitudes
as normalized inverse projection Jacobians of the stable and unstable
curves.
\end{corollary}
\begin{proof}
On the dispersing branch $\mathcal J_n=-W_{n,uv}$, and both theorems
use its exact value at the normal orbit as denominator. Uniqueness of
the limit identifies the two products. Evaluation at $v=0$ identifies
the left factors, and evaluation at $u=0$ identifies the right ones,
since each factor equals one at zero. Uniqueness of the action limit
and evaluation on the two coordinate axes prove the action assertion.
\end{proof}

\subsection{Scope of the smooth physical theorem}
The analytic product phenomenon, including its \emph{relative}
character, is therefore not by itself the distinction of
Theorem~\ref{thm:v8-main-relative}. The analytic comparison already
contains that phenomenon and gives its stable--unstable interpretation.
The function $\Delta$ alone is not the physical boundary invariant:
the projection charts occur essentially in
\eqref{eq:v14-projection-amplitudes}.

The half-line proof supplies a different, directly geometric set of
hypotheses and conclusions. It starts with arbitrary smooth contact
families and their curvature and separation bounds, constructs the
trace-class determinant factors without assuming normalizing charts,
and proves mixed geometric and endpoint estimates on common boxes.
It treats both parities and the physical first-hit channel. Positive
endpoint Hessians and the common Morse domain then transfer the
relative estimates through the actual residual-time integral, including
all fixed offset derivatives at zero. In the analytic comparison,
positivity of those Hessians is an additional geometric fact, and this
integration is a further argument; neither is part of normal-form
existence alone.

For compact analytic families the comparison is uniform when common
analytic charts, bounds and transversality margins are supplied.
Fixed mixed parameter derivatives can likewise be included when the
charts and $\Delta$ have the corresponding uniform parameter bounds:
differentiation of the normalized formulas introduces fixed powers of
$n$, which a strict exponential margin absorbs. This statement assumes
the parameter-dependent charts; it does not construct them for an
arbitrary $C^\infty$ geometric family. No impossibility of a smooth
normal-form approach is asserted.

The function-valued inverse uses the resulting physical law in an
additional way. It recovers the entire symmetrized energy pushforwards
by Volterra uniqueness, tests the mixed law by compatibility, and
transfers the inverse through an explicitly charged binary experiment.
The independent-contact jet inverse and its physical image are further
geometric statements. The normal-form comparison proves none of these
inverse or sampling conclusions. These are the precise components of
the full physical result, rather than an assertion that every local
relative product is new.
